%==============================================================================
%  Paper Review: Alias-Free Generative Adversarial Networks (StyleGAN3)
%  CSE 754 -- Deep Learning for Computer Vision -- Review Assignment (10 marks)
%
%  Compile:  pdflatex main -> bibtex main -> pdflatex main -> pdflatex main
%  (or press Recompile on Overleaf)
%==============================================================================

\documentclass[11pt,a4paper]{article}

\usepackage[utf8]{inputenc}
\usepackage[T1]{fontenc}
\usepackage[margin=2.4cm]{geometry}
\usepackage{graphicx}
\usepackage{booktabs}
\usepackage{amsmath}
\usepackage{amssymb}
\usepackage{caption}
\usepackage{xcolor}
\usepackage{microtype}
\usepackage{enumitem}
\usepackage{multirow}
\usepackage[numbers,sort&compress]{natbib}

\definecolor{NavyBlue}{HTML}{1D3557}
\definecolor{Accent}{HTML}{E07A1F}

\usepackage[colorlinks=true,linkcolor=NavyBlue,citecolor=NavyBlue,urlcolor=NavyBlue]{hyperref}
\usepackage{url}

\graphicspath{{figures/}{../figures/}}
\captionsetup{font=small,labelfont=bf}
\setlength{\parskip}{0.35em}

% A marks-visible section command, so the grader can see the mapping at a glance.
\newcommand{\marksec}[2]{%
  \section[#1]{#1\hfill\normalsize\normalfont\textcolor{Accent}{[#2]}}}

\begin{document}

%------------------------------------------------------------------- TITLE ---
\begin{center}
{\LARGE\bfseries Paper Review\par}
\vspace{0.4em}
{\Large Alias-Free Generative Adversarial Networks (StyleGAN3)\par}
\vspace{0.8em}
{\color{Accent}\rule{0.5\textwidth}{1pt}}
\vspace{0.9em}

\begin{tabular}{r@{\hspace{1.1em}}l}
\textbf{Reviewed paper} & T. Karras, M. Aittala, S. Laine, E. Härkönen, J. Hellsten,\\
                        & J. Lehtinen, T. Aila. \emph{Alias-Free Generative Adversarial}\\
                        & \emph{Networks}. NeurIPS 2021 (arXiv:2106.12423).\\[0.35em]
\textbf{Venue rank}     & NeurIPS --- A* (CORE).\\[0.35em]
\textbf{Reviewer}       & Shamima Hossain, ID 21141119\\[0.35em]
\textbf{Course}         & CSE 754 --- Deep Learning for Computer Vision\\[0.35em]
\textbf{Date}           & \today
\end{tabular}
\end{center}

\vspace{0.8em}
\hrule
\vspace{0.6em}

\noindent\textbf{Summary.}
StyleGAN3 makes a diagnosis and a repair. The diagnosis is that GAN generators
secretly depend on absolute pixel coordinates, so fine detail sticks to the image
grid instead of to the surfaces it belongs to; the paper traces this to
\emph{aliasing} introduced by the generator's own upsampling and nonlinearities.
The repair is to reinterpret every feature map as a bandlimited continuous
signal and to make each layer respect that interpretation. The outcome is a
generator that is translation- and rotation-equivariant to sub-pixel precision
while matching StyleGAN2's FID. This review argues that the diagnosis is the
paper's real contribution, that the evaluation is narrower than the claims, and
that the most interesting open question --- whether equivariance is even
desirable outside natural images --- is left untouched.

\vspace{0.4em}
\hrule
\vspace{0.8em}

%=========================================================== 1. PROBLEM  =====
\marksec{Problem Statement}{1 mark}

Generative adversarial networks~\citep{goodfellow2014gan} build images
hierarchically: coarse, low-resolution features are upsampled, mixed by
convolutions, and enriched with new detail at each scale. The implicit promise of
that design is that the position of a fine feature is inherited from the coarse
feature it belongs to --- move the head and the nose moves with it, and the skin
pores move with the nose.

The paper observes that this promise is not kept. In StyleGAN2~\citep{karras2020stylegan2},
much of the fine detail is instead pinned to \emph{absolute pixel coordinates}.
The authors name the symptom \textbf{``texture sticking''}: during a latent-space
interpolation, hair and fur do not travel with the head but remain fixed to the
screen, producing horizontal streaks when a pixel column is stacked over time
(the paper's Figure~1). Averaging images generated from a small latent
neighbourhood should give a uniformly blurry result; with StyleGAN2 some regions
stay unnaturally sharp, because they never moved.

The paper's central claim is that this is not a training artefact or a capacity
problem but a \textbf{signal-processing bug}. Generators can bypass the intended
hierarchy by reading positional information that leaks in through four side
channels: image borders, per-pixel noise inputs, positional encodings, and ---
the one the paper focuses on --- \textbf{aliasing}. Whenever a layer produces
frequencies above the Nyquist limit of its own sampling grid, those frequencies
fold back as low-frequency content that is tied to the grid rather than to the
image. The generator can then use this folded content as a coordinate reference.

\paragraph{Why it matters.} A still image whose texture is glued to the grid may
look perfectly good; the defect only becomes visible in motion. The problem is
therefore invisible to FID~\citep{heusel2017fid}, the field's dominant metric,
which is exactly why it persisted through several generations of state-of-the-art
models. The paper's stated goal is a generator whose sub-pixel feature positions
are inherited entirely from coarse features, making GANs usable for video and
animation rather than for isolated stills.

\paragraph{Assessment.} The problem statement is unusually strong. It identifies
a defect that the community's standard metric structurally cannot see, gives it a
falsifiable mechanistic cause, and motivates it with a demonstration that needs no
metric at all. The weakness --- developed in §\ref{sec:limits} --- is that the
motivating example is about \emph{local, hierarchical} motion (nose with head),
whereas what the paper ultimately delivers and measures is equivariance to
\emph{global} rigid transforms of the whole image.

%======================================================= 2. METHODOLOGY  =====
\marksec{Methodology and Analysis}{2 marks}

\subsection{The central idea: discrete arrays as continuous signals}

The paper's method follows from a single reframing. A feature map is not an array
of numbers but a \emph{sampled representation of a bandlimited continuous
signal}. By the Nyquist--Shannon theorem, a feature map sampled on a grid of
spacing $1/s$ can faithfully represent only frequencies below $s/2$, and the
discrete array $Z$ and its continuous counterpart $z$ are related by convolution
with an ideal interpolation filter,
\begin{equation}
z(x) = (\phi_s * Z)(x), \qquad
\phi_s(x) = \operatorname{sinc}(s x_0)\cdot\operatorname{sinc}(s x_1).
\end{equation}
Every operation in the generator must then be understood as an operation on the
\emph{continuous} signal. An operation is admissible only if (i) it commutes with
translation and rotation in the continuous domain, and (ii) its output remains
bandlimited to the grid that will store it. Any layer violating (ii) aliases, and
aliasing is what leaks coordinates.

This reframing immediately isolates the culprit. Convolutions and upsampling are
linear and easy to make well-behaved. \textbf{Pointwise nonlinearities are the
only operation that creates genuinely new frequency content} --- applying ReLU to
a continuous signal can introduce arbitrarily high frequencies. They are
therefore simultaneously the generator's only source of new detail and its only
source of aliasing. That observation is the paper's intellectual core.

\subsection{Measuring the property before fixing it}

Before changing anything the authors define two metrics, reported as PSNR in
decibels. For translation,
\begin{equation}
\text{EQ-T} = 10\log_{10}\!\left(
I_{\max}^{2} \big/
\mathbb{E}\Big[\big(g(t_x[z_0];w)_c(p) - t_x[g(z_0;w)]_c(p)\big)^{2}\Big]\right),
\end{equation}
comparing the generator's output after translating its input against the
translation of its output; EQ-R is the analogue for rotations drawn uniformly
from $[0^\circ,360^\circ)$. Higher is better. Defining the measurement first, then
deriving architecture changes against it, is good practice and makes the
subsequent ablation legible --- though §\ref{sec:limits} argues it also creates a
circularity the paper does not resolve.

\subsection{The nine changes}

The redesign proceeds through named configurations, each isolating one change
(Table~\ref{tab:configs}).

\begin{description}[leftmargin=1.6em,style=nextline,itemsep=0.25em]
\item[B --- Fourier features.] The learned $4{\times}4$ constant input is replaced
  by Fourier features, which define a spatially infinite map that can be
  translated and rotated \emph{exactly} rather than approximately. This is what
  makes the input side of the equivariance question well-posed at all.
\item[C --- No noise inputs.] StyleGAN2's per-pixel noise is removed. It is a
  literal coordinate reference: stochastic detail keyed to grid position. Removing
  it costs the model its source of fine stochastic variation --- a real sacrifice,
  discussed later.
\item[D --- Simplified generator.] Reduced mapping-network depth, and mixing and
  path-length regularisation disabled, to give the redesign a clean baseline.
\item[E --- Boundaries and upsampling.] Feature maps are given a small margin so
  the ``infinite extent'' assumption approximately holds near borders, and
  upsampling filters are improved.
\item[F --- Filtered nonlinearities.] The key change. The ideal operation is
  $\psi_s * \sigma(z)$: apply the nonlinearity in the continuous domain, then
  low-pass filter away the frequencies it created. This cannot be computed
  exactly, so it is approximated by \textbf{upsample $\to$ LReLU $\to$
  downsample}, with a $2\times$ temporary resolution increase found sufficient in
  practice.
\item[G --- Non-critical sampling.] Filter cutoffs are lowered below the Nyquist
  limit, deliberately sacrificing a sliver of representable bandwidth to suppress
  the aliasing that critical sampling leaves behind.
\item[H --- Transformed Fourier features.] A learned global transform is applied
  to the input features, letting the model place the whole synthesis frame.
\item[T --- Flexible layer specifications.] Per-layer cutoffs and sampling rates
  become free parameters chosen by an explicit schedule rather than inherited from
  StyleGAN2. Yields \textbf{StyleGAN3-T}.
\item[R --- Rotation equivariance.] Separable filters are replaced by radially
  symmetric $\mathrm{jinc}$ filters and $3{\times}3$ convolutions by $1{\times}1$,
  since neither separability nor square kernels are rotation-invariant. Yields
  \textbf{StyleGAN3-R}, with 56\,\% fewer parameters per layer.
\end{description}

\subsection{Analysis of the methodology}

\paragraph{What is strong.} The work is \emph{theory-first} in a literature that
is overwhelmingly empirical. The changes are not hyperparameter search; each is
derived from a stated requirement, and the ablation shows each one paying off in
the predicted direction. The identification of nonlinearities as the unique
aliasing source is genuinely illuminating and transfers beyond GANs: any
generative decoder with pointwise nonlinearities has the same defect.

\paragraph{Where theory and implementation separate.} The continuous-domain
argument is exact, but every realisation is an approximation, and the paper is
candid that the approximations are empirical. The ideal filter has infinite
support; a windowed Kaiser filter of size $n=6$ is used. The ideal nonlinearity
requires the continuous domain; $2\times$ upsampling is used because it ``has
been found sufficient''. The parameter ablation makes the gap visible: raising
upsampling to $m=4$ improves EQ-T from 66.65 to 74.21\,dB but costs $2.31\times$
training time. Equivariance is therefore not achieved but \emph{purchased}, along
a continuum, and the operating point is chosen by taste rather than derived.

\paragraph{An asymmetry the paper acknowledges but does not act on.} Only the
generator is treated. The discriminator remains an ordinary aliasing CNN, and it
is the discriminator that defines the training signal. If $D$ prefers to see
front teeth at particular pixel locations, $G$ is rewarded for putting them
there, no matter how equivariant $G$'s internals are. The paper notes exactly this
in its limitations --- and stops. Given that the same continuous-signal treatment
would apply to $D$ essentially unchanged, this is the most conspicuous piece of
unfinished work in the paper.

%============================================================ 3. RESULTS  =====
\marksec{Results and Outcomes}{2 marks}

\subsection{The configuration ladder}

\begin{table}[htbp]
\centering
\caption{Progressive results on unaligned FFHQ at $256^2$ (paper's Figure~3,
left). FID lower is better; EQ-T and EQ-R in dB, higher is better.}
\label{tab:configs}
\small
\begin{tabular}{llrrr}
\toprule
& \textbf{Configuration} & \textbf{FID} $\downarrow$ & \textbf{EQ-T} $\uparrow$ & \textbf{EQ-R} $\uparrow$ \\
\midrule
A & StyleGAN2                       & 5.14 & ---   & ---   \\
B & + Fourier features              & 4.79 & 16.23 & 10.81 \\
C & + No noise inputs               & 4.54 & 15.81 & 10.84 \\
D & + Simplified generator          & 5.21 & 19.47 & 10.41 \\
E & + Boundaries \& upsampling      & 6.02 & 24.62 & 10.97 \\
F & + Filtered nonlinearities       & 6.35 & \textbf{30.60} & 10.81 \\
G & + Non-critical sampling         & 4.78 & \textbf{43.90} & 10.84 \\
H & + Transformed Fourier features  & 4.64 & 45.20 & 10.61 \\
T & + Flexible layers (StyleGAN3-T) & 4.62 & \textbf{63.01} & 13.12 \\
R & + Rotation equiv.\ (StyleGAN3-R)& \textbf{4.50} & \textbf{66.65} & \textbf{40.48} \\
\bottomrule
\end{tabular}
\end{table}

Translation equivariance rises from 16.2\,dB to 66.7\,dB --- roughly a
$10^{5}$-fold reduction in equivariance error --- while FID ends slightly
\emph{better} than the StyleGAN2 baseline (4.50 vs 5.14). Rotation equivariance
is essentially untouched (${\sim}10.8$\,dB) until config R, which lifts it to
40.5\,dB; this is informative, confirming that translation and rotation
equivariance are separate properties requiring separate architectural treatment.

Two rows deserve attention as evidence of honest reporting. Configs E and F
\emph{degrade} FID (5.21 $\to$ 6.02 $\to$ 6.35) while improving EQ-T. The authors
keep them because the theory demands them, and the loss is recovered at config G.
A paper optimising for a headline number would have reordered or hidden this.

\subsection{Generalisation across datasets}

\begin{table}[htbp]
\centering
\caption{Results across six datasets (paper's Figure~5, left). ADA~\citep{karras2020ada}
used for the smaller sets.}
\label{tab:datasets}
\small
\begin{tabular}{llrrr}
\toprule
\textbf{Dataset} & \textbf{Model} & \textbf{FID} $\downarrow$ & \textbf{EQ-T} $\uparrow$ & \textbf{EQ-R} $\uparrow$ \\
\midrule
\multirow{1}{*}{FFHQ-U $1024^2$}
 & StyleGAN2      & 3.79 & 15.89 & 10.79 \\
 & StyleGAN3-T    & 3.67 & 61.69 & 13.95 \\
 & StyleGAN3-R    & 3.66 & 64.78 & 47.64 \\
\midrule
FFHQ $1024^2$
 & StyleGAN2      & \textbf{2.70} & 13.58 & 10.22 \\
 & StyleGAN3-T    & 2.79 & 61.21 & 13.82 \\
 & StyleGAN3-R    & 3.07 & 64.76 & 46.62 \\
\midrule
MetFaces $1024^2$
 & StyleGAN2      & 15.22 & 16.39 & 12.89 \\
 & StyleGAN3-T    & \textbf{15.11} & 65.23 & 16.82 \\
 & StyleGAN3-R    & 15.33 & 64.86 & 46.81 \\
\midrule
AFHQv2 $512^2$
 & StyleGAN2      & 4.62 & 13.83 & 11.50 \\
 & StyleGAN3-T    & \textbf{4.04} & 60.15 & 13.51 \\
 & StyleGAN3-R    & 4.40 & 64.89 & 40.34 \\
\midrule
Beaches $512^2$
 & StyleGAN2      & 5.03 & 15.73 & 12.69 \\
 & StyleGAN3-T    & \textbf{4.32} & 59.33 & 15.88 \\
 & StyleGAN3-R    & 4.57 & 63.66 & 37.42 \\
\bottomrule
\end{tabular}
\end{table}

The pattern is consistent: EQ-T lands in the 59--65\,dB range everywhere against
13--19\,dB for StyleGAN2, while FID moves by small amounts in both directions.
StyleGAN3 wins FID on three of the five sets shown and loses on aligned FFHQ
(2.70 $\to$ 2.79/3.07) --- unsurprising, since aligned FFHQ is the dataset where
a fixed coordinate frame is genuinely informative and coordinate dependence is
therefore an advantage rather than a bug.

\subsection{Cost}

\begin{table}[htbp]
\centering
\caption{Training cost on FFHQ $1024^2$ (paper, §4).}
\label{tab:cost}
\small
\begin{tabular}{lrrr}
\toprule
\textbf{Model} & \textbf{Parameters} & \textbf{GPU-hours} & \textbf{Overhead} \\
\midrule
StyleGAN2    & 30.0M & 1106 & --- \\
StyleGAN3-T  & 22.3M & 1576 & +42\,\% \\
StyleGAN3-R  & 15.8M & 2248 & +103\,\% \\
\bottomrule
\end{tabular}
\end{table}

Equivariance is not free. StyleGAN3-R uses \emph{half} the parameters of
StyleGAN2 yet takes twice as long to train, because the filtered nonlinearities
and radially symmetric filters are computationally heavy. The whole project
consumed 92 GPU-years and 225\,MWh.

\subsection{Ablations that carry information}

Several ablations are more informative than the headline numbers:
\begin{itemize}[leftmargin=1.4em,itemsep=0.2em]
  \item \textbf{Re-enabling noise inputs collapses equivariance}: EQ-T drops
        63.01 $\to$ 24.46\,dB. Direct confirmation that per-pixel noise is a
        coordinate channel, not innocuous stochasticity.
  \item \textbf{Path-length regularisation harms FID} (4.62 $\to$ 5.00) while
        \emph{improving} EQ-T. The authors reason that it penalises image change
        under latent motion and so actively encourages sticking, and suspect the
        equivariance gain comes from blurrier output --- a candid,
        self-undermining observation.
  \item \textbf{Lanczos filters damage rotation equivariance} (EQ-R 40.48 $\to$
        25.25) because separable filters are not radially symmetric --- precisely
        as the theory predicts.
  \item \textbf{A $p4$ G-CNN alternative}~\citep{cohen2016gcnn} reaches only
        17.07\,dB EQ-R at $2.48\times$ the training time, against 40.48\,dB for
        config R. Steerable filters~\citep{weiler2019e2cnn} were found infeasible
        at this network size.
  \item \textbf{Internal representations change qualitatively}: some feature maps
        switch from encoding magnitude to encoding \emph{phase}, which is what a
        network must do to track sub-pixel position.
\end{itemize}

\subsection{Assessment of the evidence}

Within its declared scope the evidence is thorough and honest. The scope,
however, is narrower than the abstract's closing claim that the results ``pave
the way for generative models better suited for video and animation''. No video
model is trained; no temporal metric such as FVD~\citep{unterthiner2018fvd} is
reported; the video evidence is a set of qualitative accompanying clips. The
quantitative case rests entirely on two metrics the paper itself introduces.

%========================================================= 4. LIMITATIONS =====
\marksec{Limitations and Research Gaps}{2 marks}
\label{sec:limits}

\subsection{Limitations the authors acknowledge}

\begin{enumerate}[leftmargin=1.5em,itemsep=0.3em]
\item \textbf{The discriminator is untreated.} Conceded to cause visible failures
      --- FFHQ teeth do not move correctly when the head turns.
\item \textbf{Assumptions about the data can be violated.} If the training data is
      itself aliased --- point-sampled cartoons, letterboxing, low-quality JPEG,
      retro pixel art --- then jagged, grid-locked edges are a \emph{defining}
      feature to be reproduced, and smooth sub-pixel translation directly
      contradicts fidelity to the data.
\item \textbf{Stochastic variation was removed, not replaced.} Noise inputs had to
      go; nothing consistent with hierarchical synthesis was put in their place.
\item \textbf{Only rigid motions are handled.} Scaling, anisotropic scaling and
      general deformations are listed as future work.
\item \textbf{Everything happens in sRGB.} Antialiasing should precede tone
      mapping; all GANs including this one operate after it.
\item \textbf{Attention- and token-based generators may be incompatible.} VQGAN-style
      tokenisers~\citep{esser2021vqgan} are flagged as ``possibly at odds'' with
      equivariance, and whether they can be made equivariant is left open.
\item \textbf{Cost.} 92 GPU-years and 225\,MWh for the project; $+103\,\%$
      training time for the rotation-equivariant configuration.
\end{enumerate}

\subsection{Gaps this review identifies}

\paragraph{G1. The metrics are self-validating.} EQ-T and EQ-R are introduced by
the same paper that optimises them, and no external task confirms that more
decibels means a better generator. The paper's own guidance --- that EQ-T above
${\sim}50$\,dB and EQ-R above ${\sim}40$\,dB ``look good'' --- is calibrated by
watching videos. This is a closed loop: the architecture is designed to satisfy a
property, the property is measured by a purpose-built metric, and the metric's
threshold is set by inspecting outputs of that architecture. Nothing here is
dishonest, but nothing independently anchors the scale either.

\paragraph{G2. The headline result is a non-regression, not an improvement.}
``Matches the FID of StyleGAN2'' means image quality did not get worse. The
entire positive value proposition therefore lives in video and animation, which is
exactly the setting that receives no quantitative evaluation. The paper's strongest
claim and its weakest evidence coincide.

\paragraph{G3. What was motivated is not what was delivered.} The opening example
is \emph{hierarchical} and \emph{local}: moving a head moves the nose, which moves
the pores. What the method achieves and measures is equivariance to \emph{global}
rigid transformation of the entire image. A generator can be perfectly equivariant
under global translation while still failing to move a nose correctly with a
turning head, because that is a local, non-rigid deformation. The FFHQ teeth
failure is precisely an instance of this, and attributing it solely to the
discriminator may be too generous to the method.

\paragraph{G4. Equivariance is assumed to be universally desirable.} The paper
treats coordinate dependence as pathology. But its own aligned-FFHQ result shows
StyleGAN2 winning on FID where the coordinate frame is meaningful, and its own
cartoon example concedes that grid-locked structure is sometimes correct. There is
no mechanism for a model to be equivariant \emph{where appropriate} and
coordinate-dependent elsewhere; the choice is made globally, at architecture
design time, by the practitioner.

\paragraph{G5. Medical and scientific imaging is a blind spot.} This gap follows
from G4 and is the one I find most consequential. In B-mode ultrasound --- the
modality of my course project --- image texture is \emph{speckle}: a coherent
interference pattern determined by the transducer geometry and the beam-forming
grid, not by the surfaces of the tissue. Speckle genuinely is ``stuck'' to the
imaging coordinate frame, and a physically faithful generator of ultrasound
\emph{should} keep it there while allowing anatomy to move independently.
StyleGAN3's architecture forbids exactly this. The same argument applies to CT
streak artefacts, MRI coil-sensitivity shading, and the burnt-in depth rulers and
on-screen chrome present in essentially every clinical ultrasound frame.

\paragraph{G6. No generality evidence beyond faces, animals and beaches.} Five of
the six datasets are natural images dominated by a single centred subject. Neither
a texture-dominated domain nor a structured scientific domain is tested, so the
claim of ``generally applicable, small architectural changes'' is broader than
what was shown.

%======================================================== 5. CONCLUDING  =====
\marksec{Concluding Remarks --- Suggested Improvements}{3 marks}

The reframing at this paper's heart --- treat feature maps as bandlimited
continuous signals, and every layer as an operator on them --- is correct, useful
and transferable, and I expect it to outlive the specific architecture. My
suggestions below are aimed at the gaps of §\ref{sec:limits}, ordered from
cheapest to most speculative.

\subsection*{S1. Apply the same treatment to the discriminator, and test the
prediction the paper already made}

The authors attribute the FFHQ teeth failure to a coordinate-dependent
discriminator but do not test it. The fix needs no new theory: the
filtered-nonlinearity construction (upsample $\to$ LReLU $\to$ downsample) and
the non-critical sampling schedule apply to $D$ unchanged. \textbf{Concretely},
train three models on FFHQ --- equivariant $G$ only (the paper), equivariant $D$
only, and both --- and report EQ-T alongside a targeted measurement of the teeth
artefact (landmark displacement of the dental region across a head-turn
interpolation). This is a falsifiable test of the paper's own hypothesis, costs
one additional training run per arm, and would either close the most obvious gap
or show that the explanation was wrong. Either outcome is publishable.

\subsection*{S2. Replace self-validating metrics with a downstream measurement}

To break the closed loop of G1, equivariance should be validated by a tool that
knows nothing about it. \textbf{Concretely}, generate latent-interpolation
sequences and measure temporal coherence with an \emph{independent} optical-flow
estimator such as RAFT~\citep{teed2020raft}: for a genuinely equivariant
generator, the flow field between consecutive frames should be smooth and should
integrate consistently over the sequence, whereas texture sticking produces
near-zero flow in textured regions while the underlying structure moves ---
a signature that is easy to quantify as the divergence between flow measured on
structure and on texture. Reporting FVD~\citep{unterthiner2018fvd} on generated
sequences would similarly convert the abstract's video claim into evidence.

I weight this suggestion heavily because of an experience in my own project: I
validated a derived orientation target using a structure tensor, then found the
number only trustworthy once I checked it against an entirely independent
principal-axis estimator. Self-consistent measurements are comfortable and can be
confidently wrong.

\subsection*{S3. Make band-limiting adaptive rather than uniform}

The paper applies the same expensive treatment at every layer, and its own
parameter ablation shows a clear Pareto frontier that it does not exploit: filter
size $n=4$ gives 57.5\,dB at $0.84\times$ time, $n=6$ gives 66.7\,dB at
$1.00\times$, upsampling $m=4$ gives 74.2\,dB at $2.31\times$. Since aliasing is
created only at nonlinearities and only where the signal actually contains
near-Nyquist energy, the filtering effort should follow the spectrum.
\textbf{Concretely}, measure per-layer spectral occupancy during training and set
each layer's upsampling factor and filter support from it, spending $m=4$ only in
the few layers that need it. Given StyleGAN3-R's $+103\,\%$ overhead, recovering
even half of that at equal EQ-R would materially change the method's
practicality.

\subsection*{S4. Target local deformation equivariance, not only global rigid motion}

This addresses G3, the gap between what the paper motivates and what it delivers.
Global translation and rotation are the trivial members of the transformation
family the introduction actually describes. \textbf{Concretely}, extend the input
Fourier-feature grid to accept a smooth, low-frequency \emph{deformation field}
rather than a rigid transform, and define EQ-D by analogy with EQ-T using
small-amplitude random warps. The continuous-signal argument carries over
directly, since a smooth warp of a bandlimited signal is still bandlimited for
sufficiently small amplitude; the bandwidth budget simply has to account for the
warp's Jacobian. This would let one state and test the paper's motivating claim
--- that a nose moves with a head --- rather than a proxy for it. Spatial
transformer machinery~\citep{jaderberg2015stn} offers a ready parameterisation.

\subsection*{S5. Selective equivariance: let the model choose its own frame}

Gaps G4 and G5 share a root cause --- equivariance is imposed globally when it is
really a property of \emph{content}. Some image content belongs to object
surfaces and should move with them; other content belongs to the imaging process
and should not. \textbf{Concretely}, I propose a \emph{two-stream generator}: an
equivariant structural stream built exactly as StyleGAN3 prescribes, and a
deliberately coordinate-dependent stream that retains per-pixel noise and
unfiltered positional encoding, with a learned spatial gate combining them. The
gate is supervised by nothing --- it is free to discover which content is
surface-bound.

This also restores what config C threw away. StyleGAN3 removed noise inputs
because they leak coordinates; the two-stream design keeps them where leaking
coordinates is physically correct.

The test case I would use is ultrasound synthesis. In B-mode imaging, anatomy
(aponeuroses, fascicles, muscle boundaries) is surface-bound and should be
equivariant, while speckle is beam-bound and should not be --- and the on-screen
depth ruler and probe annotations are rigidly fixed to the display frame. A
single generator currently cannot represent that split. A successful two-stream
model would be checkable objectively rather than by eye: measure whether the
synthesised speckle retains the correct first-order statistics (Rayleigh
distributed under fully developed speckle) while anatomical structures translate
correctly. Synthetic ultrasound with realistic speckle would be genuinely useful,
since annotated medical data is scarce --- my own project had 1\,048 annotated
aponeurosis frames in total~\citep{ritsche2024dltrack}.

\subsection*{S6. Buy rotation equivariance only where it is needed}

StyleGAN3-R doubles training time to make \emph{every} layer rotation-equivariant.
But the paper's own analysis holds that coarse layers determine structure while
fine layers add detail, and its config-T contribution made per-layer
specifications a free parameter. \textbf{Concretely}, build a hybrid: radially
symmetric $\mathrm{jinc}$ filters and $1{\times}1$ convolutions in the
low-resolution layers where global pose is decided, and cheaper translation-only
treatment at high resolution. If most of the 40.5\,dB EQ-R survives at closer to
StyleGAN3-T's $+42\,\%$ overhead rather than $+103\,\%$, rotation equivariance
becomes affordable by default instead of a deliberate trade.

\subsection*{Closing assessment}

This is a strong paper whose durable contribution is diagnostic rather than
architectural: it shows that a defect invisible to the field's dominant metric
had been present in every state-of-the-art generator, and it supplies the
signal-processing vocabulary to reason about such defects. The architecture that
follows is a competent, expensive, and somewhat rigid instantiation of that idea.
Its principal unexamined assumption --- that dependence on image coordinates is
always pathological --- happens to hold for photographs of faces and animals, and
to fail for a large class of scientific and medical imagery where the coordinate
frame is physically meaningful. The most valuable extension is therefore not more
equivariance but \emph{selective} equivariance: a generator that learns which of
its content belongs to the world and which belongs to the camera.

%------------------------------------------------------------- REFERENCES ----
\bibliographystyle{plainnat}
\bibliography{references}

\end{document}
